\subsection*{缓存租期怎样约束新写入}

前面的任务租约借助时间安排接管，资源代次检查排除迟到旧写。要判断时间本身能承担多少保护，可以回到 Gray 与 Cheriton 的 Leases~\cite{leases}。该论文研究分布式文件缓存：客户端取得数据及有效租期，在租期内使用缓存，服务器要修改被覆盖的数据，必须等各持有者同意并使本地副本失效，或者等其租期到期。

取客户端缓存值 \(v_0\)、租期 10 秒。5 秒时，服务器想写 \(v_1\)，客户端不可达，服务器继续等待。租期过去后，服务器可以按规则发布新值。客户端再次读取时，先向服务器取得有效租期并检查数据。旧副本读取和新写入在同一份数据上接受协调：旧读能够使用缓存的期间，服务器承担等待义务，等待结束后，旧客户端通过重新取得租期来恢复读取资格。

服务器崩溃重启时也要尊重先前授出的租期。论文给出等待最大租期的简单办法，或在持久存储中保留更精确的租期记录~\cite{leases}。若服务器重启即认定全部租期不存在，立即发布 \(v_1\)，原客户端可能在自认为有效的期间继续使用 \(v_0\)。

这里的时钟条件直接参与一致性。服务器时钟过快，会在客户端认为租期有效时提前允许写。客户端时钟过慢，会持续使用服务器已经视为过期的副本。论文以已知有界的时钟漂移作为最低时间条件，并讨论同步时钟与误差余量~\cite{leases}。时钟漂移指经过同一段真实时间时，两个时钟累计时长的差异。绝对时刻相差多少与计时速度相差多少，需要分别判断。

\CEightFigure{ch08-leases-wait}{缓存有效期内由服务器等待旧租期结束，发布以后客户端重新取得读取资格}{fig:ch08-leases-wait}

把这套约束放回 T1，就能看清本章租约实际保护的位置。Leases 的服务器掌握被缓存数据的新写入口，客户端也在每次使用缓存时遵守有效租期。T1 的数据库则只掌握恢复资格，按 C→T 行锁顺序串行决定谁领取、谁能改调度。W1 发到支付平台的请求已经离开这个仲裁位置，平台既看不到 T1 的到期变化，也没有参与等待旧领取结束。于是我们在本地用代次检查保护调度，在独立资源上用原子水位检查保护内容，在平台上用固定支付身份合流资金动作。每条写入路径都有实际接收者执行检查，租期才与具体效果相接。

\subsection*{Chubby 怎样把领取变化带到资源}

Chubby 面向粗粒度协调，客户端可以持锁数小时，借锁选出服务主节点。它的复制服务保存锁与少量元数据，客户端会话通过周期性的 KeepAlive 续约请求维持~\cite{chubby}。会话有自己的租期，客户端维护保守的本地到期判断，失去有效判断时停止使用缓存。这里的 Chubby 会话与购买请求身份里的 E1 分属不同协议。

锁服务可以先后把锁授予 W1、W2，W1 先前发送的请求却可能晚于 W2 的动作到达远端。Chubby 因此提供 sequencer，内容包含锁名、取得模式和锁授予代次 generation。客户端将它带到受保护服务，接收者确认是否有效、模式是否合适，再处理请求~\cite{chubby}。本章的 T1 与代次也在资源接收处核验。Chubby sequencer 还包含共享或独占模式，接收者根据所执行的动作判断所需模式。

原文给出两类校验依据：维护 Chubby 会话与一致缓存，或者只与服务器最近见过的 sequencer 比较~\cite{chubby}。后一种对应上一节的水位历史，资源未见到新代次时依据旧水位判断。采用这种方案的资源可以锁住对应记录，核对锁名、所需模式和 generation，旧代次当场拒绝，通过则在同一资源事务中推进水位、写入内容并提交，具体原子边界与第 8.4 节相同。

前一种通过与 Chubby 协调取得有效性信息，接收服务把缓存有效性及其失效协议，与本地受保护写入衔接起来。若读出有效后释放本地保护、稍后无条件写内容，就会重新留下检查后的交错位置。这份衔接由接收服务实现，Chubby 自身只提供锁与有效性信息。

对尚未支持 sequencer 的服务，Chubby 提供 \code{lock-delay}：持有者失联导致锁释放后，暂缓其他客户端取得锁，给迟到消息留下消退机会。论文明确把它称为不完善的缓解机制~\cite{chubby}。若旧消息延迟超过等待时间，它可以在新持有者写入后到达。资源检查能够在到达时拒绝它，或者系统在有可证明延迟上界的模型中，把等待时长设到覆盖该上界，二者各自提供具体的排除依据。

\subsection*{从单个 Redis 键开始}

Redis 的锁模式先解决同一时刻谁能取得一条记录，以及旧持有者怎样安全释放。设资源名为 \code{report:C1}，其锁键为 \code{lock:report:C1}，初始不存在。W1 为这次申请生成随机值 \(r_1\)，执行：

\begin{BookCode}
SET lock:report:C1 r1 NX PX 30000
\end{BookCode}

\code{NX} 要求键尚不存在，\code{PX 30000} 设置 30000 毫秒的生存时长，这个时长称为 TTL（time to live）。Redis 在同一条命令中判断键并保存值和期限，命令成功返回后，W1 才取得这一节点的领取。键存在期间，W2 对同一键的 \code{SET NX} 失败。键到期后，W2 可以用自己的随机值 \(r_2\) 重新取得它~\cite{redis-redlock}。

随机值标记一次申请，取值在不同客户端、不同申请之间具有足够低的碰撞概率。它首先用于释放时核对归属。假设 W1 暂停超过 30 秒，原键到期，W2 已写入 \(r_2\)。W1 恢复后若直接 \code{DEL}，就会删除 W2 的键。释放入口把比较与删除放在同一个原子操作中，只有当前值等于 \(r_1\) 才删除。当前 Redis 官方文档给出以下命令，并说明它从 Redis 8.4 引入~\cite{redis-redlock}：

\begin{BookCode}
DELEX lock:report:C1 IFEQ r1
\end{BookCode}

旧版本可以用原子执行的脚本完成同一项条件删除。若在客户端先 GET、再 DEL，两步之间 W2 可以创建新键，误删问题又会出现。核值必须覆盖实际删除动作。比较失败时 W1 结束释放，当前的 \(r_2\) 保留。

\CEightFigure{ch08-redis-single}{随机值使旧释放只删除自己的记录，异步复制却可能丢失已确认领取}{fig:ch08-redis-single}

随机值安全释放保护的是 Redis 里的键。W1 长暂停后向外部资源发出旧写，资源能否接受，取决于资源自己的入口。只给 Redis 命令增加核值，远端 \code{report:C1} 的内容更新尚无新判断依据。

再把单节点改成主从故障切换。主节点 M 确认 W1 的 \code{SET}，记录尚未异步复制到副本 N，M 随即退出。N 被提升为主节点，因本地没有这条键而接受 W2。W1 和 W2 各持有一次成功确认。Redis 官方用这条历史说明异步复制故障切换留下的互斥窗口~\cite{redis-redlock}。后来的删除核值只避免互相误删，先前丢失的领取历史已经影响了授予决定。

\subsection*{Redlock 怎样把拒绝依据分散到五个节点}

Redlock 采用 \(N\) 个相互独立的 Redis 主节点，节点之间不依靠复制同步这一把锁。以 \(N=5\) 为例，分别记为 A、B、C、D、E。客户端以同一键名、同一随机值向各节点申请，只有至少三个节点成功，且总耗时符合有效期要求，才把本轮视为成功~\cite{redis-redlock}。独立节点使一次节点故障可以由其他节点提供记录，客户端负责收集并判断这一组局部结果。

完整取得过程从本地采样时刻 \(t_0\) 开始。客户端向各节点发送带 TTL 的条件设置，每个节点调用使用远小于 TTL 的超时，避免一个不可达节点耗尽整轮时间。当前官方描述建议并行发送，以压缩取得过程。收集结果后再次采样 \(t_1\)，计算 \(\Delta=t_1-t_0\)，同时检查成功节点数与剩余有效期。未取得多数，或者共同窗口已经耗尽，则按随机值核对释放所有节点，包括那些响应丢失、客户端尚不确定是否成功的节点。

为什么五个节点中要有三个？假设 W1 的成功节点集合为 \(Q_1\)，W2 的为 \(Q_2\)，两者都至少有三个节点。由于总共只有五个节点，
\[
 |Q_1\cap Q_2|
 =|Q_1|+|Q_2|-|Q_1\cup Q_2|
 \geq 3+3-5=1.
\]
因此两轮成功至少共享一个节点。若 W1 的旧键在这段共同有效期间保持存在，共享节点上的 \code{NX} 就会拒绝 W2。W2 最多取得剩余两个节点，达不到多数。这是从局部拒绝推到全局互斥的关键一步。

其中有一项时间限定：交集节点必须在所讨论的期间保留有效旧记录。取 W1 在 A、B、C 成功，D、E 不可达。C 的时钟提前向前跳，使旧键过早到期，或者 C 在记录落盘前崩溃并立即空状态重启。W2 随后在 C、D、E 成功。A、B 保留 W1，C、D、E 保存 W2，两个客户端都可能认为自己的窗口尚未结束~\cite{kleppmann-locking}。

\CEightFigure{ch08-original08}{多数集合的交集只有保留有效旧记录时才能继续仲裁}{fig:ch08-quorum}

集合相交这个算式始终成立，失效的是从交集推出拒绝的前提。故障模型中若允许节点在有效期内遗忘已确认记录，就要解释它怎样恢复这份拒绝依据。给每个节点加一个副本，只要提升时可以丢失同一条记录，就会把单主故障切换的窗口带回这一层。

Redis 官方讨论了两条具体路径~\cite{redis-redlock}。一条是持久化，逐次同步使节点发出成功确认以前保住记录。只每秒同步一次时，断电可能丢失最近的键，有序关闭恢复正常也无法排除这条断电历史。另一条是延迟重新加入：可能遗忘记录的节点重启后，至少等待超过系统采用的最大 TTL，等旧领取全部失效，再参与授予。这条路径用等待换取对旧记录失效的确定判断，因而依赖可靠度量该时长。多数节点同时退出时，等待会直接限制整个服务的可用性。

\subsection*{共同有效窗口怎样计算}

五个节点接到请求的时刻各异，各键也在不同时刻开始计时。假设 W1 最早的一条键可能在 \(t_0\) 就已设置，最晚一条在取得过程结束以前设置。忽略误差时，到 \(t_1\) 判断成功的那一刻，最早的键最多已消耗 \(\Delta\)，其剩余时间至少为 \(\mathrm{TTL}-\Delta\)。其他已成功的键开始得更晚，所以也至少覆盖这一段剩余时间。

再留出时钟漂移等造成的安全余量 \(\varepsilon\)，得到所有已成功键共同存在的一段保守窗口：
\[
 V=\mathrm{TTL}-\Delta-\varepsilon.
\]
例如 \(\mathrm{TTL}=10\) 秒、\(\Delta=0.2\) 秒、\(\varepsilon=0.1\) 秒，逐项扣除后
\[
 V=10-0.2-0.1=9.7\text{ 秒}.
\]
这里 \(\Delta\) 来自客户端对整个取得过程的计时，\(\varepsilon\) 要覆盖模型允许的误差。最终取得要求成功节点达到多数，并且 \(V>0\)。即使 \(\Delta<\mathrm{TTL}\)，扣除余量后也可能已经没有窗口。它们服务于共同窗口的下界，取一个很小的经验数值，就只能支持相应的经验判断~\cite{redis-redlock}。

\CEightFigure{ch08-validity}{最早设置的键决定保守共同窗口，取得耗时与误差都从 TTL 中扣除}{fig:ch08-validity}

这项推导依赖各时钟以近似一致的速率度量经过时长。客户端与 Redis 日期显示相差几分钟，和其中一台把 10 秒计成 2 秒，是两种问题。前者是绝对时刻偏差，后者改变键究竟存在多久，会破坏共同窗口的下界。Redis 当前文档还说明 TTL 到期机制使用非单调墙上时钟，向前校时可能使键提前到期~\cite{redis-redlock}。墙上时钟提供日期与时刻，单调计时则用于度量经过时长，避免日期校正带来的回退。实现采用哪种计时、允许多大漂移和跳变，要与上述下界一起判断。

有了 9.7 秒共同窗口，客户端也只取得了这段期间内的工作机会。若用租期本身保护一个没有接收端检查的动作，相关效果要在窗口结束以前完成。工作者执行了 1 秒就发出请求，网络再延迟 20 秒，平台收到时已经越过窗口。把 TTL 增加到一分钟，只把反例中的暂停或延迟改成更长时间，允许任意长停顿的模型依旧包含这条历史。

\subsection*{两处长暂停为什么得出不同判断}

Kleppmann 检查任意长进程暂停、网络延迟、时钟异常和崩溃遗忘下的资源正确性，建议以递增 fencing token 让接收者拒绝已经落后的请求~\cite{kleppmann-locking}。antirez 的回应采用有界相对计时误差，强调取得多数后再次检查总耗时，把取得过程中的延迟算进 \(\Delta\)~\cite{antirez-redlock}。把两处暂停分别放进历史，才能看清这些判断作用于哪一步。

第一条历史中，W1 在等待节点响应时暂停 20 秒，随后恢复，收到了原先的成功响应。若 \(\mathrm{TTL}=10\) 秒，取得后的再次采样会得到 \(\Delta>20\) 秒，剩余窗口为负，这次获取失败，W1 按核值释放局部记录。antirez 指出的后置耗时检查能够处理这个位置的暂停。用早已接收到的三个成功响应直接宣布取得，就遗漏了算法规定的检查。

第二条历史中，W1 已完成多数判断，测得 \(\Delta=0.2\) 秒，计算出 9.7 秒窗口，随即在发送资源写入以前暂停 20 秒。W2 可以在旧键到期后取得多数并写入新内容。W1 恢复后继续旧代码，资源若无接收端核验，就可能接受旧写。后置耗时检查在暂停以前已经结束。即使把检查挪到发送前，暂停也可以发生在最后一次检查与生效之间。

\CEightFigure{ch08-pauses}{取得期间的暂停进入耗时检查，最后检查后的暂停由资源入口处理}{fig:ch08-pauses}

antirez 的回应也把最后检查后的暂停归入自动过期锁的使用问题。这一点同时适用于本章数据库租约、Chubby 会话和 Redis 键：授予服务能按自己的规则更新持有者，远端资源要根据实际到达的请求决定是否生效。允许任意长暂停时，继续写入的权限就要在能控制效果的位置兑现。

\subsection*{随机值的相等性与代次的先后关系}

antirez 还提出用随机 token 配合资源的检查并设置~\cite{antirez-redlock}。设资源已经把当前值从 \(r_1\) 原子安装为 \(r_2\)，写入要求请求 token 等于当前值，W1 带 \(r_1\) 的旧写就被拒绝。这项保护来自资源维护当前资格和原子比较，随机值用来回答请求是否属于同一次安装。

安装规则本身同样要处理迟到。若任何声称持锁的请求都可以无条件安装自己的值，W1 在 W2 之后把 \(r_1\) 装回去，就重新获得了表面匹配。可行协议要为安装规定合法前驱、当前资格或共同仲裁，让已经撤销的持有者无法凭迟到请求覆盖当前值。这样设计时，资源的比较与安装已经承担了一部分协调责任。

图~\ref{fig:ch08-random-install} 取其中一种安装规则：资源初始保存 \(r_1\)，W1 暂停后，W2 的安装同时核验合法前驱与当前资格，原子改为 \(r_2\)。W1 随后发出的旧写与迟到安装分别进入自己的核验入口。

\CEightFigure{ch08-random-install}{随机值的相等比较保护当前写入，合法安装规则阻止旧值迟到后重新取得资格}{fig:ch08-random-install}

递增 token 提供了另一种可利用的信息。接收者见过 8 后，单凭 \(7<8\) 就能拒绝迟到安装或写入，水位只向前移动。随机 \(r_1,r_2\) 之间只有相等与不同，数值大小未必对应授予先后。两种办法都要把检查与业务写入原子衔接，递增水位还要保证分配顺序、作用域和重启后的历史。只保存最大已观察值时，其保护起点仍是第 8.4 节所说明的资源观察时刻。

这也解释了两篇文章采用的不同评判位置。Kleppmann 要求在长暂停、迟到及时钟异常模型中保护最终资源，故把检查责任放到接收者。antirez 先讨论有界计时条件下的授予，再讨论随机 token 和资源比较能怎样保护使用。这些结论要连同各自假设阅读。Redis 当前官方文档保留共同窗口和相对时钟速率的前提，同时在一致性提示中建议使用 fencing token，并提醒进程活着和锁仍有效之间可能分离~\cite{redis-redlock}。

\subsection*{恢复机会怎样继续出现}

TTL 让失联客户端留下的键最终到期，失败申请者主动释放局部键，可以更早给其他客户端机会。多个客户端反复竞争时，可能各占少数节点，谁都达不到三个。官方建议失败后随机延迟再试，减少持续同步碰撞~\cite{redis-redlock}。随机退避改善取得机会，确定的单次成功依旧经过多数和有效期判断。

网络分区使客户端无法主动释放，其他客户端可能等待 TTL。若分区和竞争一直持续，等待可以反复发生，多个节点可用这一事实本身尚未给出某位客户端必然取得的时刻。续期则在当前有效窗口内，对同一随机值条件延长各节点期限，在原窗口结束前完成多数续期，并扣除本轮耗时与误差重新计算有效窗口，才把延长视为成功。失联或缺少成功确认时，客户端根据原到期边界停止使用。一个已经过期的客户端去重新争取记录，是新一轮竞争，要根据当前记录重新判断。无限延长也可能使其他参与者始终没有机会，因此延长次数、调度和竞争处理共同影响进展。

对 P1 而言，增加 Redis 锁可以减少几个扫描器同时准备同一项工作的成本，但恢复任务与支付事实已经位于购买库，当前 C→T 事务直接决定调度资格。若在外层加 Redis，工作者取得它后仍进入这个数据库决定点，最后执行的本地修改由当前 T 约束。原平台键 \code{pay-C1} 则处理可能重叠的支付到达。这样分配责任时，外层协调的偶发重复会增加竞争或调用次数，资金效果的上限由固定支付身份及其保护范围承担。

\subsection*{回到 P1 的安全与进展}

现在把各层接到同一次购买。A 保存 \(S\)、P1、C1 和 T1。调度器在数据库中领取 T1，第 7 代失联以后第 8 代可接管。所有可能生效的支付到达使用原平台范围、\code{pay-C1} 和原输入。可信结果经过 C1 合流，R 首次建立 O1 与稳定结果，并结束 T1。旧调度写入在当前代次处停止，旧工作者持有的可信证据则经过支付事实核验进入 R。

安全关系逐项落在对应迁移上。领取、续期与接管只改变恢复资格，保留 P1 的业务阶段及固定支付映射。R 依据支付证据推进终态与任务完成，重复证据恢复原 O1。资源端在同一原子边界检查代次与写入，保持水位不下降，进度覆盖再核对业务版本。平台效果上限由第七章的身份保护范围确定。租期结束只改变接管条件，最终拒绝证据同时确认原身份既往无成功，并排除已受理、在途及未来尝试成功。

取 T1 已由 W2 接管第 8 代，W1 随后带回真实 X1，平台余额为 7000，与此同时资源只保存第 7 代水位。W1 修改 T1 调度会因旧代次被拒绝，其可信 X1 可以进入 R，资源在尚未接受 8 以前则可以按水位 7 接受第 7 代写。这三个结果分别来自任务、支付证据与资源的实际观察，单看 W1 是否为当前持有者无法决定全部结果。

持续进展取决于实际的执行机会。每项到期任务被持续扫描并最终获得调度，候选之间公平处理，下次运行时间是有限且可达到的时刻。数据库恢复可用，持锁事务最终结束，时钟推进到接管条件，新领取者有足够连续运行时间完成必要步骤。若每轮领取者都在工作以前退出，任务可以一直存在，付款却始终没有执行机会。

对尚未发生 E 的 C1，自动执行还要求平台和网络在安全发送窗口内提供受理机会，所有可能生效的到达符合原身份保护规则，并由平台处理到可解释结论。任意长故障若跨过固定窗口，任务转为核查，后续接管者继续查询原付款或取得账单证据。

对 E 已发生的 C1，平台结果保留且能凭原映射取得，至少一条可信结果通道最终可用，R 获得本地提交机会。任务进入核查后，核查者取得足以确定结局的证据，再提交对应结论。平台永久丢失可关联证据时，本地继续如实保存未决，完成所需的证据机会已经缺失。Alice 最后获知结果，则依赖一次查询或响应完整到达并被保存。

P1 现已能够跨进程持续恢复，并把平台结果补齐到唯一 O1。购买服务接下来要通知履约系统：若 R 已经提交，发送通知的进程随即退出，履约系统可能完全不知道这份已付款订单。通知本身需要保留身份、待办事实与新的恢复路径。
